\section*{Nomenclatura}

\begin{tabbing}
$A$\qquad\qquad \= Área da seção transversal (m\textsuperscript{2})\\
$B$ \> Largura da pista (m), equivalente a $b_{w,\mathrm{pista}}$\\
$\ell_t$ \> Vão de cálculo da prancha do tabuleiro (m)\\
$b_{tr}$ \> Largura tributária de uma longarina (m)\\
$a$ \> Espaçamento entre eixos do veículo (m)\\
$P$ \> Carga característica por roda (kN)\\
$b_w$ \> Largura da seção transversal da peça do tabuleiro (m)\\
$b_{w,\mathrm{pista}}$ \> Largura da pista de rolamento (m)\\
$CIV$ \> Coeficiente de impacto vertical\\
$d$ \> Diâmetro da longarina circular de madeira roliça (m)\\
$esp_{long}$ \> Espaçamento entre longarinas (m)\\
$esp_{tab}$ \> Espaçamento entre peças do tabuleiro (m)\\
$E_{c0,med}$ \> Módulo de elasticidade médio na compressão paralela às fibras (MPa)\\
$E_{M0}$ \> Módulo de flexão estimado na modelagem, não medido na fonte (MPa)\\
$E_{0,\mathrm{m}}$ \> Módulo de elasticidade médio à flexão adotado no dimensionamento (MPa)\\
$E_{t0}$ \> Módulo de elasticidade na tração paralela às fibras (MPa)\\
$\alpha_E$ \> Fator adotado na relação $E_{M0}=\alpha_E E_{c0}$\\
$f_{c0}$ \> Resistência à compressão paralela às fibras (MPa)\\
$f_{m}$ \> Resistência à flexão (MPa)\\
$f_{v0}$ \> Resistência ao cisalhamento paralelo às fibras (MPa)\\
$f_m(\mathbf{x})$ \> $m$-ésima função objetivo\\
$g_j(\mathbf{x})$ \> $j$-ésima restrição de desigualdade\\
$h$ \> Altura da seção transversal da peça do tabuleiro (m)\\
$I$ \> Momento de inércia da seção (m\textsuperscript{4})\\
$k_{mod}$ \> Coeficiente de modificação\\
$L$ \> Vão teórico (m)\\
$N_c$ \> Número de realizações perturbadas na avaliação robusta\\
$V_{\mathrm{total}}$ \> Volume total de madeira da superestrutura (m\textsuperscript{3})\\
$W$ \> Módulo de resistência da seção (m\textsuperscript{3})\\
$\mathbf{x}$ \> Vetor de variáveis de projeto\\
$\gamma_g,\ \gamma_q$ \> Coeficientes parciais das ações permanentes e variáveis\\
$\gamma_{wf},\ \gamma_{wc}$ \> Coeficientes parciais do material (tensão normal e cisalhamento)\\
$\delta_{\mathrm{tot}}$ \> Flecha final quase permanente no meio do vão (m)\\
$\Delta V$ \> Diferença relativa de volume entre os cenários C e R (\%)\\
$\varepsilon_E$ \> Erro de rigidez, $100(E_{c0,c_s}/E_{c0,s}-1)$ (\%)\\
$\varepsilon_\delta$ \> Erro relativo de previsão de deslocamento (\%)\\
$U_{fin},U_Q$ \> Utilizações da flecha final quase permanente e variável\\
$U_\delta^{C\rightarrow R}$ \> Utilização final do projeto C reavaliado com rigidez R\\
C, R \> Cenários de rigidez de classe e experimental, respectivamente\\
$\varphi$ \> Coeficiente de fluência\\
$\psi_2$ \> Fator de combinação quase permanente\\
$\rho$ \> Nível de perturbação das variáveis de projeto (robustez)\\
$\rho_{ap}$ \> Massa específica aparente a 12\% de umidade (kg/m\textsuperscript{3})\\
$S_{T_i}$ \> Índice de Sobol de ordem total da variável $i$\\
Índices \> $m$: valor médio; $k$: valor característico; $d$: valor de cálculo
\end{tabbing}
